\documentclass[10pt,conference,compsocconf]{IEEEtran}
\usepackage{style}
\usepackage{tgpagella}


\begin{document}
\title{
    A Comparative Analysis of Optimization Trajectories and Convergence Properties.
}
\author{
    Ozair Faizan \\
        \texttt{\href{mailto:ahmad.faizan@epfl.ch}{ahmad.faizan@epfl.ch}} \and
    Paul Tercier \\
        \texttt{\href{mailto:paul.tercier@epfl.ch}{paul.tercier@epfl.ch}}\\
}
\maketitle

\setlength{\belowdisplayskip}{2pt} \setlength{\belowdisplayshortskip}{2pt}
\setlength{\abovedisplayskip}{2pt} \setlength{\abovedisplayshortskip}{2pt}


\begin{abstract}
    The choice of the optimization algorithm determines training efficiency and generalization of
    deep neural networks. While methods such as Stochastic Gradient Descent and Adam have become
    field standards, recent works such as Lion \cite{chen2025lionsecretlysolvesconstrained}
    and Muon \cite{jordan2024muon} promise better convergence properties in practice.
    In this paper we present a comparative study of the optimization trajectories produced
    by Stochastic Gradient Descent, Adam, Lion and Muon applied to Multi-Layer
    Perceptron on the Fashion-MNIST dataset \cite{xiao2017/online}.
    To this end, we use a linear interpolation method: for each discrete update
    step from $\theta_t$ to $\theta_{t+1}$, we evaluate the loss function at $N=100?$ points along
    the connecting vector, allowing us to empirically verify convexity and smoothness properties
    along the trajectory.

    Furthermore, we evaluate the convergence points $\vec{m}$ by testing the condition:
    $$ f(\vec{m}) \leqslant f(\vec{x}) + \nabla f(\vec{m})\T (\vec{m} - \vec{x}), $$
    for points $\vec{x}$ within a neighborhood $B_\varepsilon(\vec{m})$. 

    We use constant learning rate found by grid-search and compare them with predicted theoretical
    values, all other hyperparameters are set to default.
\end{abstract}

\section{Introduction}

\section{Models and Methods}

\section{Results}

\section{Discussion}

\section{Summary}


\newpage
\bibliographystyle{IEEEtran}
\bibliography{literature}

\end{document}


